\chapter{Introduction}
\label{ch:introduction}

\section{Background and Motivation}
Intersections are the critical bottlenecks and primary hazard zones of urban road transportation, accounting for over 40\% of vehicular collisions and 25\% of traffic fatalities. Traditional fixed-time or sensor-actuated traffic lights introduce mandatory stop-and-go idling cycles, substantially elevating fuel consumption, emissions, and commuter delays. With Autonomous Electric Vehicles (AEVs) and Connected and Automated Vehicles (CAVs), Autonomous Intersection Management (AIM) enables vehicles to negotiate right-of-way cooperatively without physical signals. However, unsignalized intersection negotiation is characterized by severe non-linearities: vehicles must navigate multi-directional conflicting trajectories (left turns, straight crossings, right turns) amidst surrounding human drivers exhibiting stochastic headway gaps. Deep Reinforcement Learning (DRL) provides an expressive framework for synthesizing safety-critical driving policies directly from high-dimensional observations with microsecond-scale online execution.

\section{Complex Engineering Problem (CEP) Justification}
\label{sec:cep}
% Strict constraint: Max 70 words justifying why this is a Complex Engineering Problem
This project addresses an unsignalized intersection negotiation problem exhibiting non-linear multi-agent stochasticity, uncoordinated human driver behavior, and dynamic occlusions. The solution requires formulating a continuous-state, discrete-action Partially Observable Markov Decision Process (POMDP) that reconciles strictly conflicting engineering objectives: collision-free safety versus velocity-maximizing transit efficiency under microsecond real-time execution constraints. These multi-faceted technical requirements, absence of analytical closed-form solutions, and strict real-time constraints rigorously classify this problem as a Complex Engineering Problem under Washington Accord criteria.

\section{Problem Statement and Research Objectives}
The central challenge is synthesizing a robust longitudinal policy for an autonomous ego vehicle navigating a four-way unsignalized intersection across three stochastic tasks: Left Turn, Straight Crossing, and Right Turn. The agent must traverse collision-free while minimizing transit time. Four concrete objectives are pursued: (i)~\textbf{Literature Taxonomy}: Establish a 5-axis taxonomy of autonomous intersection management; (ii)~\textbf{PPO Benchmark}: Implement and benchmark an on-policy PPO baseline over 1.2M episodes; (iii)~\textbf{Discrete SAC Reproduction}: Execute a 1:1 mathematical reproduction of Xiao et al. \cite{xiao2024decision} (twin critics, auto-tuned temperature $\alpha$) over 50k episodes per config; and (iv)~\textbf{Hardware Optimization \& Telemetry}: Pin CPU threads to physical P-cores to eliminate OpenMP barrier stalls, and conduct a forensic safety audit of Sim-to-Sim domain gaps.

\section{Research Challenges and Main Contributions}
Negotiation entails asymmetric geometry (7 conflict points for left turns vs. 6 straight, 2 right), sparse rewards, and microscopic yielding disparities. This project delivers five key contributions: (1)~\textbf{5-Axis AIM Taxonomy} classifying control, architecture, scheduling, flow modeling, and autonomy mixture; (2)~\textbf{PPO Benchmark} evaluating 1.2M rollouts comparing 60-dim baseline against 61-dim intent states in Ray RLlib; (3)~\textbf{1:1 Discrete SAC Reproduction} implementing the exact formulation of Xiao et al. \cite{xiao2024decision} over 100,000 cumulative episodes; (4)~\textbf{P-Core Affinity Pinning} achieving a \textbf{62\% runtime reduction} (saving 22.4 hours); and (5)~\textbf{Forensic Safety Audit of the 99.52\% Metric} analyzing high-momentum crossing dynamics under $\gamma=0.90$, microscopic IDM yielding, and the 5 recorded edge-case collisions.

\section{Report Organization}
Chapter \ref{ch:related_works} reviews literature and the 5-axis taxonomy; Chapter \ref{ch:problem_scope} formalizes geometry and POMDP; Chapter \ref{ch:methodology} presents methodology and algorithm selection; Chapter \ref{ch:implementation} details simulation, SAC, and P-core optimization; Chapter \ref{ch:results} provides results and the forensic audit; Chapter \ref{ch:ethics} discusses ethics; Chapter \ref{ch:limitations} evaluates limitations; and Chapter \ref{ch:conclusion} presents conclusions.


